\section{Counting determinant-zero triples}\label{zero:section}

The remaining part of the weighted count concerns collinear projected
points.  We organize these triples by their integral linear relation.
The auxiliary modulus excludes short relations when the three residues
are affinely independent.  When two residues coincide, the equality
usually imposes an additional lattice condition; the relations for
which it does not are sufficiently sparse.

Keep the digit matrix $C$ fixed, and write $\Lambda=\Lambda_C$ for
its row lattice from Lemma~\ref{orb:diagonal}.  Write
\[
 h=B^k,\qquad D=B^b,\qquad E=B^e,\qquad I=DE,
 \qquad 0\le b\le e\le k.
\]
Thus $C$ has diagonal form $(1,D,E)$ over $\mathbb Z/h\mathbb Z$,
$[\mathbb Z^3:\Lambda]=I$, and $E\mathbb Z^3\subseteq\Lambda$.
Recall also $H=h^2$, $q>h^{100}$, the box
$\mathcal B=[0,N)^2\times[N,2N)\cap\mathbb Z^3$, and the weight
$W(A)$ from \eqref{aux:weight}.  For a column $u\in\mathcal B$, its
projection means $(u_1/u_3,u_2/u_3)$.

\begin{proposition}\label{zero:weighted-zero}
For the lattice $\Lambda$ and weight $W$ above, let $\mathcal Z$ be
the set of integer $3\times3$ matrices whose columns belong to
$\mathcal B$, whose rows belong to $\Lambda$, whose determinant is
zero, and whose three projected columns are pairwise distinct.  Then
\begin{equation}\label{zero:target}
 \sum_{A\in\mathcal Z}W(A)
 \ll \log(2N)\frac{N^6}{I^2}.
\end{equation}
The implied constant is independent of the fixed digit matrix $C$.
\end{proposition}

The orbit restriction in Proposition~\ref{sam:weighted-count} only
reduces this sum.  We first record the lattice estimates needed for
the null vectors.

\begin{lemma}\label{zero:plane}
For a primitive vector $x\in\mathbb Z^3$, put
\[
 \Lambda_x=\Lambda\cap x^\perp,\qquad
 x\cdot\Lambda=g(x)\mathbb Z,\qquad
 J_x=\det(\Lambda_x),
\]
where $g(x)>0$.  Then
\begin{equation}\label{zero:plane-data}
 g(x)\mid E,\qquad J_x=\frac{I|x|}{g(x)},\qquad
 \frac{g(x)^3}{I}\le h^2.
\end{equation}
The reduction of $\Lambda_x$ modulo $q$ is the entire plane
\[
 \Pi_x=\{v\in\mathbb F_q^3:v\cdot\bar x=0\},
 \qquad \bar x=x\bmod q\ne0.
\]
If $\Gamma\subseteq\Lambda_x$ is a sublattice of index $m$, the
number of ordered triples of vectors in $\Gamma$, all of length at
most $4N$ and spanning a plane over $\mathbb Q$, is at most
\begin{equation}\label{zero:row-count}
 \ll \frac{N^6}{(mJ_x)^3}.
\end{equation}
\end{lemma}

\begin{proof}
Primitivity gives $x\cdot\mathbb Z^3=\mathbb Z$.  Since
$E\mathbb Z^3\subseteq\Lambda$, we have
$E\mathbb Z\subseteq g(x)\mathbb Z$, so $g(x)\mid E$.
Lemma~\ref{lat:plane-covolume} gives the formula for $J_x$, and
$g(x)^3/I\le E^2/D\le h^2$.

Primitivity also implies $\bar x\ne0$, with no restriction on the
size of $x$.  To prove surjectivity, choose $z\in\mathbb Z^3$ with
$x\cdot z=1$.  If an integral lift $v$ of a point of $\Pi_x$ has
$x\cdot v=qa$, then $v-qaz\in\mathbb Z^3\cap x^\perp$ is another
lift of the same point.  Thus $\mathbb Z^3\cap x^\perp$ reduces
onto $\Pi_x$.  Its multiple by $E$ lies in $\Lambda_x$ and still
reduces onto $\Pi_x$, because $q\nmid E$.

Finally, a triple counted in \eqref{zero:row-count} contains two
independent vectors of length at most $4N$.  Therefore the second
successive minimum of $\Gamma$ is at most $4N$ whenever the count
is nonzero.  Lemma~\ref{lat:plane-count} then bounds the choices for
each vector by $O(N^2/\det\Gamma)$, and
$\det\Gamma=mJ_x$.
\end{proof}

\begin{lemma}\label{zero:moment}
For $R\ge h$,
\begin{equation}\label{zero:moment-bound}
 \sum_{\substack{x\in\mathbb Z^3\ \mathrm{primitive}\\
                  R\le |x|<2R}}g(x)^3
 \ll R^3 I.
\end{equation}
In particular, this estimate holds when $R>H/2$.
\end{lemma}

\begin{proof}
Since $g(x)\mid B^e$, write $g(x)=B^{a(x)}$ with
$0\le a(x)\le e$.  For $0\le j\le e$, the condition
$B^j\mid g(x)$ is equivalent to
$Cx=0\pmod{B^j}$.  Indeed, if $\widetilde C$ is any integral lift
of $C$, then $\Lambda$ is the sum of the integral row span of
$\widetilde C$ and $h\mathbb Z^3$, and $B^j\mid h$.
The diagonal form of $C$ shows that the fraction of residue
classes modulo $B^j$ satisfying this congruence is
\[
 B^{-j-\max(j-b,0)}.
\]
Indeed, the first diagonal coordinate must vanish, the second
must be divisible by $B^{\max(j-b,0)}$, and the third is
unrestricted since $j\le e$.

Because $R\ge h\ge B^j$, each residue class contains
$O((R/B^j)^3)$ integral vectors of length less than $2R$.
We may discard primitivity for this upper bound.  Consequently,
\begin{align*}
 \sum_{\substack{x\ \mathrm{primitive}\\R\le |x|<2R}}g(x)^3
 &\le \sum_{j=0}^e B^{3j}
       \#\{x\in\mathbb Z^3:|x|<2R,\ Cx=0\pmod{B^j}\}\\
 &\ll R^3\sum_{j=0}^e B^{2j-\max(j-b,0)}
 \ll R^3 B^{b+e}=R^3I.
\end{align*}
The exponents in the last sum strictly increase to $b+e$, so its
bound is uniform in $B$.  Finally, $h\ge2$ implies
$H/2=h^2/2\ge h$.
\end{proof}

\begin{proof}[Proof of Proposition~\ref{zero:weighted-zero}]
Every $A\in\mathcal Z$ has rational rank two.  Indeed, its columns
are nonzero, and no two are proportional because their projections
are distinct.  Its right nullspace therefore has a primitive
integral generator $x$, unique up to sign.  Every coordinate of
$x$ is nonzero: otherwise the relation would make two columns
proportional.  Taking the cross product of two independent rows
and dividing by the gcd of its coordinates gives
\begin{equation}\label{zero:null-range}
 Ax=0,\qquad x_1x_2x_3\ne0,\qquad |x|\ll N^2.
\end{equation}
Every row of $A$ lies in $\Lambda_x$ and has length at most $4N$.
Thus $x$ records a relation among the columns, while the rows are lattice
vectors perpendicular to $x$.  We may sum over both choices of sign of
$x$ for an upper bound.

Partition the possible vectors into dyadic shells
$R\le |x|<2R$, $R=1,2,4,\ldots$.  By
\eqref{zero:null-range}, there are $O(\log(2N))$ shells.  We will
bound the weighted contribution of each shell by $O(N^6/I^2)$.
For a fixed primitive $x$ in such a shell, a condition confining every
row to an index-$m$ sublattice of $\Lambda_x$ gives at most
\begin{equation}\label{zero:fixed-null-count}
 \ll \frac{N^6g(x)^3}{m^3I^3R^3}
\end{equation}
matrices of rational rank two, by \eqref{zero:row-count} and
$J_x=I|x|/g(x)$.  In the affine case we will sum $g(x)^3$ using
Lemma~\ref{zero:moment}.  In the equal-pair cases we instead use
$g(x)^3\le Ih^2$, together with an extra row congruence or a smaller
number of possible null vectors.
The support statement and the three weight bounds in
Proposition~\ref{aux:weights} give the following exhaustive cases.

\paragraph{Affinely independent residues.}
If $W(A)>0$ and the columns of $A\bmod q$ are affinely
independent, Lemma~\ref{aux:random-set} excludes the
relation $Ax=0$ when $|x|\le H$.  To see both parts of this
exclusion, if $x_1+x_2+x_3\ne0$ use the restricted shift; if the
sum is zero, use affine independence and $\bar x\ne0$.
Hence a contributing shell satisfies $R>H/2$.
Using $W(A)\ll1$, \eqref{zero:fixed-null-count} with $m=1$, and
Lemma~\ref{zero:moment}, its contribution is at most
\begin{equation}\label{zero:affine-count}
 \ll \frac{N^6}{I^3R^3}
       \sum_{\substack{x\ \mathrm{primitive}\\R\le |x|<2R}}g(x)^3
 \ll \frac{N^6}{I^2}.
\end{equation}

\paragraph{An equal pair with an additional equation.}
Every other matrix of positive weight has two equal columns
modulo $q$.  Fix one pair $i<j$ with this property, and put
$\delta=\mathbf e_i-\mathbf e_j$, where the $\mathbf e_i$ are the
standard basis vectors.  Summing over the three pairs only
changes the implied constant.  First suppose that
$\bar x\notin\mathbb F_q\delta$.
The functional $v\mapsto v\cdot\delta$ is nonzero on $\Pi_x$:
the annihilator of $\Pi_x$ is precisely $\mathbb F_q\bar x$.
By Lemma~\ref{zero:plane}, the rows of $A$ therefore belong to
the index-$q$ sublattice
\[
 \{v\in\Lambda_x:v_i=v_j\pmod q\}.
\]
There are $O(R^3)$ integral $x$ in any shell.  Using
\eqref{zero:fixed-null-count} with $m=q$, the pointwise bound
$g(x)^3\le Ih^2$, and $W(A)\ll(q^3/s)^2$, the contribution is at most
\begin{equation}\label{zero:equal-additional}
 \ll \frac{N^6}{I^2}\,h^2q^{-3}(q^3/s)^2
 \ll \frac{N^6}{I^2}\,\frac{h^{26}}q
 \ll \frac{N^6}{I^2}.
\end{equation}
Here and below we use $q^3/s\ll qH^6=qh^{12}$ and $q>h^{100}$.
This argument includes matrices of every rank modulo $q$.

\paragraph{An equal pair without an additional equation.}
It remains to consider $\bar x\in\mathbb F_q\delta$.
The coordinate of $x$ outside the pair $i,j$ is divisible by
$q$ and, by \eqref{zero:null-range}, is a nonzero integer.
Thus $|x|\ge q$, and a contributing shell has $R>q/2$.
The nonzero multiples of $\delta$ give $q-1$ residue classes
for $x$ modulo $q$.  Each class contains
$O((1+R/q)^3)$ vectors of length less than $2R$, so the number
of possible $x$ in this shell is
\begin{equation}\label{zero:special-normals}
 \ll R^3/q^2.
\end{equation}
\smallskip
\noindent\textbf{Rank at least two modulo $q$.}
The weight bound
$W(A)\ll q^3/s$ now gives
\begin{equation}\label{zero:special-rank-two}
 \ll \frac{N^6}{I^2}\,h^2q^{-2}(q^3/s)
 \ll \frac{N^6}{I^2}\,\frac{h^{14}}q
 \ll \frac{N^6}{I^2}
\end{equation}
for the contribution of the shell.

\smallskip
\noindent\textbf{Rank at most one modulo $q$.}
All rows reduce into
one of the $q+1$ one-dimensional linear subspaces of $\Pi_x$.
This also covers the zero row space.  For each such line
$\ell\subseteq\Pi_x$, surjectivity in Lemma~\ref{zero:plane}
shows that
\[
 \{v\in\Lambda_x:v\bmod q\in\ell\}
\]
has index $q$ in $\Lambda_x$.  Formula~\eqref{zero:row-count}
thus bounds the rank-two rational row triples, summed over
all these lines, by
\[
 \ll(q+1)\frac{N^6}{(qJ_x)^3}
 \ll\frac{N^6}{q^2J_x^3}.
\]
Using \eqref{zero:special-normals} and the general weight bound,
the weighted contribution is at most
\begin{equation}\label{zero:special-rank-one}
 \ll \frac{N^6}{I^2}\,h^2q^{-4}(q^3/s)^2
 \ll \frac{N^6}{I^2}\,\frac{h^{26}}{q^2}
 \ll \frac{N^6}{I^2}.
\end{equation}

The affine and equal-pair cases cover the support of $W$.
Their bounds are uniform over the dyadic shells, so summing
over the $O(\log(2N))$ shells proves \eqref{zero:target}.
\end{proof}
